\documentclass[12pt]{article}

\usepackage[a4paper,margin=1in]{geometry}
\usepackage{amsmath,amssymb}
\usepackage{newtxtext,newtxmath}
\usepackage{fancyhdr}
\usepackage{listings}
\usepackage{graphicx}
\usepackage{float}
\usepackage{booktabs}
\usepackage{xcolor}
\usepackage{hyperref}

\setlength{\parindent}{0pt}
\setlength{\parskip}{5pt}

\AtBeginDocument{
    \fontsize{12pt}{14pt}\selectfont
}

% Header
\pagestyle{fancy}
\fancyhf{}
\fancyhead[L]{Name: \underline{Shaik Jnaved}}
\fancyhead[R]{Enrollment No.: \underline{P26DS001}}
\renewcommand{\headrulewidth}{0.4pt}
\setlength{\headheight}{15pt}

% Code formatting
\lstset{
    language=C,
    basicstyle=\ttfamily\fontsize{8pt}{10pt}\selectfont,
    breaklines=true,
    frame=single,
    columns=fullflexible,
    keepspaces=true,
    showstringspaces=false,
    keywordstyle=\color{blue}\textbf,
    commentstyle=\color{green!50!black}\textit,
    tabsize=4,
    numbers=left,
    numberstyle=\tiny\color{gray},
    numbersep=6pt
}

\begin{document}

% ==============================================================================
% PROBLEM 1: GENETIC ALGORITHM
% ==============================================================================

\begin{center}
\textbf{Computer Science and Engineering Department, SVNIT, Surat}\\
\textbf{M.Tech. I -- Semester I}\\
\textbf{High Performance Computing}\\
\textbf{Lab Assignment 5}
\end{center}

\vspace{0.2cm}

\textbf{\large Problem Statement 1}

\textbf{Title:} Task Allocation in Distributed Systems using Parallel Genetic Algorithm (OpenMP).

\textbf{Objective:}
Implement a Parallel Genetic Algorithm (GA) using OpenMP to solve the Task Allocation Problem (TAP) in distributed systems. Consider $M = 300$ tasks and $P = 10$ processing nodes ($P \ll M$). Map each task to an appropriate node such that the Total Cost (Execution Cost + Inter-task Communication Cost) is minimized. Inter-task communication cost between two tasks is incurred if and only if they are scheduled on different nodes ($a_i \ne a_j$). Compare the parallel OpenMP solution with the sequential baseline.

\vspace{0.2cm}

\textbf{Source Code (\texttt{lab\_5/main.c}):}
\lstinputlisting{../../lab_5/main.c}

\newpage

\textbf{Output (Genetic Algorithm):}
\begin{figure}[H]
\centering
\includegraphics[width=0.92\textwidth]{images/ga_output.png}
\caption{Benchmark Execution Output for Sequential vs. Parallel Genetic Algorithm ($M = 300, P = 10$)}
\end{figure}

\textbf{Analysis \& Discussion (Genetic Algorithm):}
\begin{enumerate}
    \item \textbf{Cost Model:} Total Cost comprises Execution Cost ($\sum E[i][a_i]$) and Communication Cost ($\sum_{i < j} C[i][j] \cdot [a_i \ne a_j]$).
    \item \textbf{OpenMP Parallelization:} Fitness evaluation and offspring reproduction are distributed across threads using \texttt{\#pragma omp parallel for schedule(static)}.
    \item \textbf{Thread Safety:} A lightweight thread-safe generator (\texttt{simple\_rand}) with private seeds avoids mutex contention from standard \texttt{rand()}.
    \item \textbf{Performance:} Sequential GA took \textbf{0.2810 s} vs. Parallel GA \textbf{0.0580 s}, achieving a \textbf{4.84$\times$ speedup} (30.28\% parallel efficiency on 16 threads).
\end{enumerate}

\newpage

% ==============================================================================
% PROBLEM 2: PARTICLE SWARM OPTIMIZATION
% ==============================================================================

\begin{center}
\textbf{Computer Science and Engineering Department, SVNIT, Surat}\\
\textbf{M.Tech. I -- Semester I}\\
\textbf{High Performance Computing}\\
\textbf{Lab Assignment 5}
\end{center}

\vspace{0.2cm}

\textbf{\large Problem Statement 2}

\textbf{Title:} Task Allocation in Distributed Systems using Parallel Particle Swarm Optimization (OpenMP).

\textbf{Objective:}
Implement a Parallel Particle Swarm Optimization (PSO) algorithm using OpenMP to solve the Task Allocation Problem (TAP). Represent particle positions in continuous space mapped to discrete node assignments using floor projection ($\lfloor pos[k] \rfloor$). Parallelize particle trajectory updates and fitness evaluations with OpenMP and compare execution runtime with the sequential implementation.

\vspace{0.2cm}

\textbf{Source Code (\texttt{lab\_5/pso.c}):}
\lstinputlisting{../../lab_5/pso.c}

\newpage

\textbf{Output (Particle Swarm Optimization):}
\begin{figure}[H]
\centering
\includegraphics[width=0.92\textwidth]{images/pso_output.png}
\caption{Benchmark Execution Output for Sequential vs. Parallel PSO ($M = 300, P = 10$)}
\end{figure}

\textbf{Analysis \& Discussion (Particle Swarm Optimization):}
\begin{enumerate}
    \item \textbf{Continuous-to-Discrete Mapping:} Continuous particle positions in $[0, P)^M$ are discretized via $\lfloor pos[k] \rfloor$ to compute allocation costs while maintaining smooth velocity updates.
    \item \textbf{OpenMP Parallelization:} Particle updates and objective evaluations are parallelized using \texttt{\#pragma omp parallel for schedule(static)}.
    \item \textbf{Performance:} Sequential PSO took \textbf{0.3160 s} vs. Parallel PSO \textbf{0.0450 s}, achieving a \textbf{7.02$\times$ speedup} (43.89\% parallel efficiency on 16 threads).
\end{enumerate}

\vspace{0.3cm}

\textbf{\large Summary Comparison Table}

\begin{table}[H]
\centering
\begin{tabular}{lcccc}
\toprule
\textbf{Algorithm} & \textbf{Sequential Time} & \textbf{Parallel Time (OMP)} & \textbf{Speedup} & \textbf{Parallel Efficiency} \\
\midrule
Genetic Algorithm (GA) & 0.2810 s & 0.0580 s & \textbf{4.84$\times$} & 30.28\% \\
Particle Swarm Optimization (PSO) & 0.3160 s & 0.0450 s & \textbf{7.02$\times$} & 43.89\% \\
\bottomrule
\end{tabular}
\end{table}

\end{document}
